\documentclass[12pt]{article}

% -------------------------------------------------
%                 PACKAGES
% -------------------------------------------------
\usepackage[margin=1in]{geometry}
\usepackage{setspace}
\usepackage{amsmath, amssymb, amsfonts}
\usepackage{graphicx}
\usepackage{caption, subcaption}
\usepackage{booktabs, array}
\usepackage{natbib}
\usepackage{hyperref}
\usepackage{xcolor}
\usepackage{footmisc}
\usepackage{pdflscape}
\usepackage{appendix}
\usepackage{todonotes}
\usepackage{footmisc}
\usepackage[normalem]{ulem} % provides \sout
\usepackage{xcolor}
\usepackage[draft]{changes}

\definecolor{addcolor}{HTML}{1B7F3B}
\definecolor{delcolor}{HTML}{B22222}
\definecolor{replcolor}{HTML}{4B4BA5}

\setaddedmarkup{\textcolor{addcolor}{#1}}
\setdeletedmarkup{\textcolor{delcolor}{\sout{#1}}}
\makeatletter
\renewcommand{\replaced}[2]{%
  {\textcolor{replcolor}{#1}}{\textcolor{delcolor}{\sout{#2}}}%
}
\makeatother
% -------------------------------------------------
%               SPACING
% -------------------------------------------------
\onehalfspacing

% -------------------------------------------------
%            THEOREM ENVIRONMENTS
% -------------------------------------------------
\newtheorem{theorem}{Theorem}
\newtheorem{proposition}{Proposition}
\newtheorem{corollary}{Corollary}
\newtheorem{hypothesis}{Hypothesis}

\newenvironment{proof}[1][Proof]{%
\noindent\textbf{#1.} }{\hfill$\square$}

% -------------------------------------------------
%              CUSTOM COMMANDS
% -------------------------------------------------
\newcommand{\red}[1]{\textcolor{red}{#1}}
\newcommand{\blue}[1]{\textcolor{blue}{#1}}

\newcolumntype{L}[1]{>{\raggedright\arraybackslash}m{#1}}
\newcolumntype{C}[1]{>{\centering\arraybackslash}m{#1}}
\newcolumntype{R}[1]{>{\raggedleft\arraybackslash}m{#1}}

% -------------------------------------------------
%                 DOCUMENT
% -------------------------------------------------
\begin{document}

% -------------------------------------------------
%                  TITLE
% -------------------------------------------------
\title{ECON490 Project :D\thanks{Acknowledgment text here.}}
\author{
Erika Salvador\thanks{Amherst College, esalvador28@amherst.edu} 
\and 
Caroline Theoharides\thanks{Amherst College, ctheoharides@amherst.edu.}
\and
Susan Godlonton\thanks{Williams College, sg5@williams.edu}
}
\date{\today}
\maketitle

% -------------------------------------------------
%                 ABSTRACT
% -------------------------------------------------
\begin{abstract}
Placeholder abstract text.

\medskip
\noindent\textbf{Keywords:} key1, key2, key3 \\
\textbf{JEL Codes:} key1, key2, key3
\end{abstract}

\thispagestyle{empty}
\clearpage
\pagenumbering{arabic}

% -------------------------------------------------
%               MAIN TEXT
% -------------------------------------------------

\section{Introduction}
\label{sec:introduction}

\section{Background}
\label{sec:background}

International migration is a defining feature of economic life in many low- and middle-income countries. In the Philippines, overseas employment has been actively promoted as a development strategy for decades, and remittances constitute a substantial share of national income. A large empirical literature documents that migration is associated with higher household consumption and better risk-sharing arrangements, as well as greater investment in children’s human capital \citep{yang2008international, mckenzie2007network}. These findings have shaped the prevailing view that migration improves welfare in origin communities through income gains and risk diversification.

Less understood, however, is how migration reshapes household organization and the distribution of authority within marriage. When one spouse migrates, prolonged physical separation alters both the allocation of income and the process through which financial decisions are made. \replaced
{Earnings generated abroad may remain under the migrant’s control, who determines the timing and amount of remittances, or they may be pooled and managed by the spouse who remains at home, who then exercises discretion over household spending decisions.}{Earnings generated abroad may be retained and transferred at the migrant’s discretion, or they may be pooled and administered by the spouse who remains at home.} The timing, amount, and earmarking of remittances therefore become instruments of control. Because bargaining models of the household assign decision weight to partners based on their control over resources and outside options, such shifts in financial authority have direct implications for intrahousehold power. In standard collective and noncooperative models, improvements in women’s income or fallback position increase their bargaining weight and expand their influence over allocation outcomes \citep{manser1980marriage, mcelroy1981nash}. Extensions that incorporate violence further imply that stronger economic positions for women can reduce abuse by increasing the expected cost to partners of engaging in violent behavior or by strengthening women’s capacity to exit abusive relationships \citep{aizer2010gender}.

At the same time, alternative theoretical perspectives emphasize that changes in relative economic status may generate conflict rather than cooperation. When increases in women’s income challenge established gender hierarchies, men may perceive a loss of status or authority within the household \citep{eswaran2011domestic, anderberg2016unemployment}. In such settings, shifts in economic control can alter expectations about decision-making, spending priorities, or autonomy, potentially creating tension over household roles. \added{Although women's labor force participation and international migration rates are high in the Philippines, ethnographic \citep{medina2001filipino} evidence suggests that expectations surrounding male provision and authority within marriage remain influential.} \replaced{When marital roles assign primary financial responsibility to men,}{If male identity is closely tied to the role of primary provider} a decline in relative earnings or authority may provoke attempts to reassert control, including through coercive or violent behavior. The empirical literature reflects this ambiguity. Positive labor market shocks to women are associated with reductions in intimate partner violence (IPV) in some contexts, whereas adverse shocks to men are associated with increases in abuse \citep{aizer2010gender, autor2019work}. These findings indicate that the effect of economic change on household conflict depends on how income shocks interact with bargaining positions and prevailing gender norms.

This theoretical ambiguity is likely to be especially pronounced in settings where formal exit options are limited. The Philippines provides such an environment. The country is among the world’s largest exporters of labor, with women comprising a substantial share of deployed overseas workers \citep{mcauliffe2024world}, while legal divorce remains unavailable.\footnote{Annulment is legally available but requires judicial proceedings, entails substantial financial and evidentiary burdens, and remains inaccessible for many households.} Formal marital dissolution is therefore not a credible outside option for most couples, including those experiencing IPV. In standard bargaining frameworks, the credibility of exit determines bargaining weight. When legal exit is constrained, alternative margins—such as physical separation or changes in control over income—may partially substitute for formal dissolution. International migration thus introduces prolonged spousal separation without terminating the marital contract. \added{On one hand, migration} may strengthen women’s autonomy by altering resource control and by increasing geographic distance from a partner. \added{On the other hand,} separation without dissolution may generate instability by weakening monitoring and increasing uncertainty about authority within marriage. Whether migration reduces or intensifies conflict in this setting is therefore theoretically indeterminate. \added{Migration may reduce violence by improving effective outside options, or it may intensify conflict by disrupting established household roles.} The welfare consequences of migration depend not only on income gains but also on how migration interacts with constrained exit options to reshape exposure to violence within marriage.

\added{At the same time, the relationship may also operate in reverse. Households experiencing conflict or abuse may have different incentives to send a spouse abroad. Migration may function as a means of physical separation, income stabilization, or informal exit in the presence of violence. Observed associations between migration and intimate partner violence may therefore reflect not only the effect of migration on household conflict, but also the effect of violence on migration decisions. Identifying the causal relationship requires variation that is plausibly independent of household-level conflict and allows these channels to be separated.}

\added{One source of such variation arises from differences in local enforcement conditions. Because abusive behavior is shaped not only by the expected consequences imposed by the legal environment, variation in enforcement provides a natural setting in which to study these effects. In economic models of crime, behavior depends on the expected probability and severity of punishment \citep{becker1968crime}. Extensions of this framework to domestic violence similarly highlight that abuse reflects both private bargaining dynamics and the expected cost imposed by legal sanctions \citep{aizer2010gender, iyengar2009does}. We therefore consider two distinct mechanisms: one concerning direct deterrence and the other concerning moderation. First, expanded local reporting and protection infrastructure may directly reduce violence by increasing the expected cost of abusive behavior. If victims can more easily report abuse and access formal protection, the likelihood that violence results in sanctions may arise, which strengthens deterrence independent of changes in household bargaining positions. Second, enforcement may shape how economic shocks affect household conflict. If international migration alters bargaining positions within marriage—by shifting control over income, geographic separation, or effective outside options—its impact on violence may depend on whether abusive behavior carries meaningful expected consequences. Where enforcement is weak, changes in bargaining power may not translate into reduced violence. Conversely, where enforcement is credible, shifts in economic control may more effectively constrain abusive behavior. In this case, observed violence reflects both the direct deterrent effect of enforcement and its interaction with underlying economic conditions.}

\added{In the Philippines, such variation arrises from the staggered roll out of Violence Against Women and Children (VAWC) desks across municipalities..... will continue writing}


\deleted{In the Philippines, enforcement capacity and access to protective mechanisms may plausibly vary across localities, depending on the credibility of sanctions and the protection available to victims. Even if migration strengthens women’s bargaining position within marriage, its effect on violence may depend on whether abusive behavior carries meaningful expected consequences. Where enforcement is limited, shifts in income control or outside options may have attenuated deterrent effects. Where enforcement is credible, improvements in women’s economic position may translate more directly into reductions in abuse. The impact of migration on intimate partner violence therefore depends on the interaction between household bargaining incentives and local enforcement capacity.}

Accordingly, this paper examines how international migration affects intimate partner violence in the Philippines and whether this relationship depends on local enforcement capacity. We study whether exposure to migration shifts the equilibrium incidence of violence in origin communities and whether access to formal reporting and protection mechanisms moderates this effect. To address these questions, we exploit the staggered rollout of Violence Against Women and Children (VAWC) desks across municipalities. Mandated under the Magna Carta of Women and related implementing legislation beginning in 2009, VAWC desks expanded local reporting and referral infrastructure at different points in time across jurisdictions. This staggered institutional expansion provides plausibly exogenous variation in enforcement capacity. Leveraging this variation allows us to examine whether migration and local enforcement interact in shaping observed violence outcomes and to distinguish changes in deterrence from shifts attributable solely to economic bargaining dynamics.

By integrating international migration with variation in enforcement capacity, this paper contributes to two strands of literature. First, it builds on work demonstrating that enforcement institutions shape gender-based violence. \citet{sviatschi2024gender} shows that the expansion of women’s justice centers in Peru reduces violence and improves long-term outcomes of women, which illustrates how institutional infrastructure can alter equilibrium behavior within households. We extend this line of inquiry by introducing a large, externally driven economic shock and examining how its effects depend on local enforcement capacity. Rather than estimating the average impact of institutional expansion, we also analyze how exposure to global labor market opportunities conditions the effectiveness of enforcement in deterring violence. Second, the paper expands the migration literature by incorporating intimate partner violence into the welfare analysis of migration. Existing studies emphasize consumption, education, and labor supply responses. Yet if migration reshapes bargaining positions in a context where formal exit is constrained, its welfare implications depend on how marital relationships adjust in equilibrium. The interaction between migration exposure and enforcement capacity allows the paper to characterize how labor mobility jointly influences household welfare and the incidence of violence.

More broadly, this analysis advances a framework in which household outcomes are jointly determined by economic shocks and enforcement capacity. Migration reallocates resources and reshapes bargaining positions within marriage, while enforcement capacity governs the expected consequences of abusive behavior. The equilibrium incidence of violence reflects the interaction of these forces. Migration may strengthen women’s effective outside options where enforcement is credible, but it may generate instability where sanctions are weak. A full evaluation of migration policy therefore requires accounting not only for income gains but also for the enforcement conditions that mediate its effects on household conflict and women’s safety.

\section{Literature Review}
\label{sec:literature}

\section{Data}
\label{sec:data}

\section{Methodology}
\label{sec:method}

\section{Results}
\label{sec:results}

\section{Discussion}
\label{sec:discussion}

\section{Conclusion}
\label{sec:conclusion}

% -------------------------------------------------
%               REFERENCES
% -------------------------------------------------

\clearpage
\singlespacing
\setlength{\bibsep}{0pt}

\bibliographystyle{apalike}
\bibliography{references}   

% -------------------------------------------------
%            TABLES AND FIGURES
% -------------------------------------------------

\clearpage
\onehalfspacing

\section*{Tables}
\addcontentsline{toc}{section}{Tables}

% Example:
% \begin{table}[htbp]
% \centering
% \caption{Placeholder}
% \begin{tabular}{lcc}
% \toprule
% Variable & Model 1 & Model 2 \\
% \midrule
% X & 0.123 & 0.456 \\
% \bottomrule
% \end{tabular}
% \end{table}

\clearpage

\section*{Figures}
\addcontentsline{toc}{section}{Figures}

% Example:
% \begin{figure}[htbp]
%   \centering
%   \includegraphics[width=.6\textwidth]{fig/placeholder.pdf}
%   \caption{Placeholder}
% \end{figure}

% -------------------------------------------------
%                 APPENDIX
% -------------------------------------------------

\clearpage
\begin{appendices}
\onehalfspacing

\section{Additional Results}
\label{appendix:A}

\section{Data Construction}
\label{appendix:B}

\end{appendices}

\end{document}
